\chapter{GAME DESIGN}

\section{Game Name}

The title of the game is [GAME NAME HERE]. The name reflects the game's overarching narrative, which centers on a neighborhood chemist helping community members with everyday challenges. It also incorporates wordplay, as the term "bubbly" refers both to the character's personality and to the game's bubble-shooter mechanics.

\section{Game Overview}

[GAME NAME HERE] is a 2D bubble-shooter educational game designed to teach Grade 6 chemistry concepts through interactive gameplay. Players create and transform materials using five action orbs, each representing a specific physical or chemical process. In addition to bubble-shooter levels, players can explore a village environment and assist characters through a sorting minigame focused on identifying salvageable and unsalvageable materials.

\section{Game-Concept}

[GAME NAME HERE] is a 2D stylized educational game set in a brightly colored neighborhood and follows a chemist as she assists the community with daily tasks. The challenges focus on core chemistry topics such as states of matter and, more importantly, distinguishing between physical and chemical changes in everyday life. In early levels, particularly those set in a café, players use action orbs on ingredients to observe changes in properties and achieve goals defined by a recipe system.

Progress is communicated through a traffic-light feedback system, where red indicates no progress, yellow indicates partial progress, and green indicates completion. Players must strategically use action orbs such as cooling, heating, mixing, slicing, crushing, and swapping to achieve correct outcomes for each level.

\section{Genre}

The game combines educational gaming with puzzle and simulation elements. The bubble-shooter mechanics place it within the puzzle genre, emphasizing logic, accuracy, and planning. The neighborhood exploration component functions as a light simulation, reinforcing learning through contextual storytelling and interaction. This hybrid genre approach supports both active recall and contextual understanding.

\section{Target Audience}

The primary target audience of [GAME NAME HERE] is Grade 6 students aligned with the MATATAG curriculum. At this stage, learners are introduced to foundational chemistry concepts such as states of matter, physical and chemical changes, and mixtures. The game's visual style, feedback systems, and mobile-first design make it accessible to a wide range of learners who benefit from interactive and gamified instruction.

\section{Game Flow Summary}

The game progresses through increasingly challenging levels, beginning in the café location. As players advance, new ingredients, action orbs, and mechanics are introduced. Bubble-shooter levels demonstrate various physical transformations such as melting, freezing, cutting, and crushing. Sorting minigame levels also increase in complexity by introducing a greater variety of materials.

\section{Look and Feel}

The game adopts a cartoon-inspired visual style influenced by titles such as Animal Crossing: New Horizons, Kirby's Dream Buffet, and Hello Kitty Island Adventure. Design inspiration is also drawn from PopMart blind box figures, aligning with the game's whimsical and approachable aesthetic. All assets are 2D, stylized, and rounded to emphasize the game's bubbly theme.

\section{How does the game insert itself in a pedagogical scenario?}

[GAME NAME HERE] integrates into pedagogical scenarios by aligning gameplay mechanics with core chemistry learning competencies. It functions as a supplementary instructional tool that reinforces classroom learning through interactive engagement. Teachers may use the game to introduce or review concepts related to physical and chemical changes without the constraints of a physical laboratory. Each level functions as a short lesson paired with a problem-solving activity, providing objectives, challenges, and immediate feedback.

Beyond the bubble-shooter mechanics, the village environment enables contextual application of learned concepts. Players interact with characters, explore locations, and encounter progressively challenging scenarios that connect chemistry concepts to familiar situations. This narrative-driven exploration promotes curiosity and reinforces understanding through storytelling.

The game is designed around two learning objectives aligned with Bloom's Taxonomy.

\textbf{LO1 (Lower-Order):} Classify physical changes and chemical changes by identifying states of matter and recognizing whether an action results in a physical or chemical transformation, including examples of uniform and non-uniform mixtures.

\textbf{LO2 (Higher-Order):} Explain and justify whether a change is physical or chemical by applying evidence from simulated real-world scenarios, such as predicting outcomes of heating or cooling materials and selecting appropriate actions to achieve desired results.

Through the Village and Bubble Shooter Lab components, [GAME NAME HERE] supports inquiry, experimentation, and critical thinking. Immediate feedback mechanisms, including character reactions and post-level explanations, help correct misconceptions and reinforce accurate understanding.

\section{Gameplay and Mechanics}

This section will be talking about the gameplay and mechanics that are going to be shown by the game [GAME NAME HERE]. This section will talk about how the game is played and how the game will supply the necessary learning objectives to the players.

\subsection{Gameplay}

Core gameplay takes place in a café setting where players complete chemistry-based challenges focused on changes in states of matter, physical and chemical changes, and the effects of heat on materials. Players launch action orbs such as Heat, Cool, Crush, and Cut toward target materials to trigger appropriate reactions. Incorrect actions result in penalties accompanied by visual and auditory feedback. As players progress, challenges increase in complexity through multi-step reactions and environmental modifiers, with guidance provided to support understanding.

\subsection{Moment to Moment}

% Game Loop Table
% TODO: Add content from original document (lines 465-488)

\begin{longtable}{|p{0.3\textwidth}|p{0.6\textwidth}|}
\hline
\textbf{Phase} & \textbf{Description} \\
\hline
\endfirsthead

\hline
\textbf{Phase} & \textbf{Description} \\
\hline
\endhead

\hline
\endfoot

% Add game loop phases here

\caption{Game Loop Table}
\label{tab:game_loop}
\end{longtable}


\subsection{Mechanics}

Set within a narrative-driven environment, players complete tasks inspired by everyday situations, such as food preparation and material sorting. Action orbs simulate scientific processes in an accessible arcade-style format. Progress is visualized through a traffic-light system and a three-star scoring mechanic, enabling players to easily assess their performance.

\subsection{Physics}

The game employs basic bubble-shooter physics to simulate predictable trajectories and collisions. These mechanics support intuitive gameplay while reinforcing timing and spatial reasoning skills.

\subsection{Movement}

[GAME NAME HERE] uses a top-down perspective for both gameplay and exploration. A virtual joystick enables movement and aiming, allowing seamless control across mobile and desktop platforms.

\subsection{Objects}

Interactive objects include action orbs, reaction tiles, collectible tokens, compound-forming elements, and village features such as character kiosks and informational displays. Each object supports learning, progression, or narrative development.

\subsection{Actions}

Player actions include aiming and shooting action orbs, switching between available orbs, interacting with characters, unlocking new areas, and viewing educational information within the village.

\subsection{Combat}

The game does not include traditional combat. Instead, challenges are presented as puzzles involving resource management, timing, and condition-based problem solving.

\subsection{Economy}

Gameplay emphasizes resource management. Players must efficiently use limited material orbs to complete level objectives, encouraging critical thinking and strategic planning.

\subsection{Game Options}

[GAME NAME HERE] includes adjustable sound and control sensitivity settings, a contextual help menu, and a pause menu. These options enhance usability and player comfort without interfering with educational objectives.

\subsection{Constraints due to the pedagogical objectives}

To maintain pedagogical integrity, gameplay mechanics are intentionally simplified and scaffolded. All reactions align with Grade 6 learning outcomes, and pacing is designed to minimize stress while reinforcing understanding.

\section{Story, Setting and Character}

\subsection{Story and Narrative}

The narrative follows a chemist assisting in the opening of a neighborhood café. Through completing recipes and tasks, the player observes how ingredients undergo physical and chemical changes. Later challenges introduce irreversible chemical changes, reinforcing key distinctions between transformation types and encouraging reflection on real-world applications.

\subsection{Game progression}

Progression is structured to introduce new concepts gradually, allowing players time to practice and internalize information before advancing. Narrative and gameplay progression are designed to support continuous learning without cognitive overload.

\subsection{License considerations}

All third-party and open-source assets will be properly licensed and credited. In-house code will follow MIT licensing, while creative assets will use Creative Commons Attribution-NonCommercial 4.0 (CC BY-NC 4.0).

\subsection{Game World}

The game world centers on a chemistry-themed café presented in a stylized 2D environment. The setting blends cozy design with experimental elements, providing a playful space for discovery, experimentation, and problem solving.

\subsection{Characters}

Players assume the role of a research assistant guided by the Local Bubble Chemist. Through completing tasks and experiments, players unlock new mechanics and increasingly complex challenges.

\subsection{Levels}

Levels follow a classic bubble-shooter structure inspired by games such as Zuma, combined with exploration and interaction reminiscent of life-simulation games. Friendly guidance and narrative context support learning progression.

\subsection{Training levels}

The training level introduces the game's mechanics and educational goals in a low-pressure environment. Visual cues, simple prompts, and immediate feedback build player confidence and prepare learners for more complex challenges.

\subsection{Assessment}

Player analytics will be integrated to track engagement and learning-related metrics such as playtime and level completion. These data will support analysis of learning progress and engagement.

\section{Interface}

\subsection{Visual System}

A top-down view ensures clear visibility of the play area. Interface elements are minimal and designed to present essential information without obstructing gameplay.

\subsection{Control System}

The game uses screen-based controls optimized for both mouse and touchscreen input, ensuring accessibility across platforms.

\subsection{Audio, music, sound effects}

Audio assets will be sourced from in-house production, open-source libraries, and licensed third-party resources, with proper attribution provided.

\subsection{Help System}

The in-game help system provides tutorials, definitions, and visual examples accessible from the main menu and pause screen, allowing learners to review concepts at their own pace.
